\documentclass[12pt]{article}
% Basic packages
\usepackage[utf8]{inputenc}   % UTF-8 encoding
\usepackage{geometry}          % Page margins
\usepackage{setspace}          % Double spacing
\usepackage{amsmath, amssymb} % Math symbols
\usepackage{graphicx}         % Include images
\usepackage{booktabs}         % Nice tables
\usepackage{enumitem}         % Custom lists
\usepackage{cancel}
\usepackage{float}

% Page setup: 1-inch margins
\geometry{
  letterpaper,
  margin=1in
}

% Double spacing
\doublespacing{}

% START OF DOCUMENT %
\begin{document}

\begin{center}
    {\textbf{Johnny A. Rojas \& Aiden Sitorus, George Nicacio, Jeremiah
        Hasibuan, Matthew Miller}}\\
    {\Large \textbf{A LAB REPORT ON KIRCHHOFF'S LAWS OR RULES}}
\end{center}

\section{Theory}
Building upon our previous analysis—where a plumbing system was used as an
analogy for electrical flow—this report evaluates the validity and application
of Kirchhoff’s Rules in complex circuits that cannot be simplified through basic
series or parallel configurations. Kirchhoff’s Rules provide the mathematical
framework necessary to solve for unknown variables in multi-loop circuits
through two fundamental principles:

\subsection*{I. Kirchhoff’s Voltage Law (KVL)}

KVL is a reflection of the Law of Conservation of Energy. It states that the
algebraic sum of the potential differences (voltages) around any closed loop or
``conductive path'' must equal zero. Much like a ``summation of forces'' in
mechanical equilibrium, KVL ensures that the energy supplied by a source is
exactly accounted for by the potential drops across the loads.

Mathematically, this is expressed as:
\begin{equation*}
  \sum_{k=1}^{n} V_k = 0
\end{equation*}
where:
\begin{itemize}
  \item $n$ represents the total number of electrical components encountered
        along a specific closed path.
  \item $k$ serves as the index for each individual element, effectively
        ``counting'' our progress as we traverse the loop.
  \item $V_k$ is the potential difference (measured in Volts) across the $k^{th}$
        element; by convention, a voltage source (battery) represents a positive
        gain, while a resistor represents a negative drop.
\end{itemize}

\subsection*{II.\@ Kirchhoff’s Current Law (KCL)}

KCL is a reflection of the Law of Conservation of Charge. It dictates that
charge cannot ``pile up'' or vanish at a junction (node). While our previous
analogy compared current to water pressure, KCL specifically describes the flow:
the total current entering a junction must exactly equal the total current
exiting it.

At any junction, the sum of currents is zero:
\begin{equation*}
    I_{in} = I_{out_1} + I_{out_2} + \cdots + I_{out_n}
  \end{equation*}
where:
\begin{itemize}
  \item $I$ represents the electric current, measured in Amperes ($A$).
  \item $n$ represents the number of outgoing branches from the junction.
\end{itemize}

To test the efficacy of these models, our team analyzed two circuits of varying
complexity. This study serves to validate the accuracy of Kirchhoff’s Rules and
to exercise advanced circuit analysis techniques, providing a foundation for the
study of more sophisticated electrical systems.

\section{Conclusion}
\subsection*{Activity 1: First Circuit}
In Activity 1 our team used multimeters to measure the experimental currents and
voltage drops across three resistors ($2220$ $\Omega$, $1494$ $\Omega$,
and $820$ $\Omega$) and compared them to theoretical values derived from a
system of linear equations. Our measurements yielded mixed results; for
the $2220$ $\Omega$ and $1494$ $\Omega$ resistors,the measured currents
of $0.0016$ A and $0.0012$ exhibited percent errors of $12.57\%$ and $7.460\%$
from the expected values of $0.0018$ A and $0.0013$A, respectively.
The $820$ $\Omega$ resistor presented a significant anomaly, with a measured
voltage drop of $1.48$ V deviating from the expected $0.4373$ V by $238.4\%$.
These discrepancies, particularly the extreme error in the third resistor,
suggest that while the circuit followed a recognizable pattern, the results were
heavily impacted by systematic errors such as internal meter resistance or
contact issues in the breadboard assembly.

\subsection*{Activity 2: Second Circuit}
For Activity 2, our team calculated the expected theoretical currents across
each node. The experimental data demonstrated varying levels of precision; we
observed relatively accurate current measurements for the $2220$ $\Omega$
and $1800$ $\Omega$ resistors, recording values of $0.00003070$ A and $0.0008$
A, which yielded low percent errors of $0.1478\%$ and $1.113\%$ compared to
theoretical calculations. Conversely, readings for
the $820$ $\Omega$, $1494$ $\Omega$, and $1200$ $\Omega$ resistors deviated more
significantly, with current errors of $58.37\%$, $27.28\%$, and $23.83\%$
respectively. Despite these fluctuations, the general pattern of the data
broadly aligns with the theoretical framework of Kirchhoff's rules. The larger
discrepancies are likely attributable to procedural and systematic inaccuracies
during data collection—such as inconsistent probe contact or minor assembly
errors—rather than a flaw in the underlying physical laws.

\section{Questions}
\subsection*{Activity 1 Questions}
\begin{enumerate}
  \item \textbf{Compare your results in step \#2 and \#3. Discuss.} \\
        After analyzing and comparing our results to the analog model introduced
        in the previous laboratory session, we observed that the values at each
        node are influenced by preceding circuit elements. This supports the
        conclusion that current flows in a ``river-like'' pattern; however, it
        is important to note that this macro-scale pattern does not necessarily
        reflect electron behavior at the microscopic level. Instead, distinct
        physical mechanisms govern charge transport to allow this observed
        behavior to occur.
\end{enumerate}

\section{Above and Beyond}
While these labs have allowed us to experience firsthand how to understand the
changes in voltage and current across a circuit, it was not until we commenced
research for this report that we discovered that
the flow of energy across wires
has less to do with the actual physical displacement of electrons and more to do
with the electric fields they produce. It is fascinating how these mechanisms
produce this predictable behavior at the macroscopic level. However, further
research also provided insight into the limitations that occur in extreme
cases—such as why certain wires require specific insulation for protection or
the challenges of long-distance intercommunication. While the majority of
circuit analysis applications can be successfully modeled using Kirchhoff's
rules, advanced applications where accuracy and timing are fundamental require
us to extend our models to account for these underlying physical complexities.

\end{document}
